% figures/information/pim_family.tex
%
% PIM Family
% Usage:
% \begin{figure}[ht]
% \centering
% \input{figures/information/pim_family}
% \caption{BIM and AIM as complementary specializations within PIM.}
% \label{fig:pim-family}
% \end{figure}

\begin{tikzpicture}[
  node distance=1.2cm,
  box/.style={
    draw,
    rounded corners,
    align=center,
    minimum width=4.2cm,
    minimum height=0.9cm,
    font=\small
  },
  note/.style={
    align=center,
    font=\footnotesize
  },
  arrow/.style={->, thick}
]

\node[box] (pim) {PIM\\Process Information Modelling};

\node[box, below left=1.4cm and 1.3cm of pim] (bim)
{BIM\\Object- and asset-oriented};

\node[box, below right=1.4cm and 1.3cm of pim] (aim)
{AIM\\Alignment- and state-oriented};

\node[note, below=0.55cm of bim] (bimnote)
{assets\\components\\lifecycle};

\node[note, below=0.55cm of aim] (aimnote)
{alignments\\states\\operations};

\draw[arrow] (pim) -- (bim);
\draw[arrow] (pim) -- (aim);

\end{tikzpicture}
